\section{Ausblick auf die Bachelorarbeit}

Aufbauend auf den Ergebnissen dieses Projekts ergeben sich viele Möglichkeiten zur Weiterentwicklung, die sich gut im Rahmen einer Bachelorarbeit umsetzen lassen.
Ein naheliegender nächster Schritt ist es, das bestehende Framework um Funktionen zur Modellierung stationärer Roboter zu erweitern. Ein Knickarmroboter würde sich dafür besonders gut eigenen.

Der Fokus dieser Erweiterung soll auf der Umsetzung der Vorwärts- und Rück\-wärts\-ki\-ne\-ma\-tik liegen. Während die Vorwärtskinematik beschreibt,
wie man aus bekannten Gelenkwinkeln die Position des Endeffektors berechnet, geht es bei der Rück\-wärts\-ki\-ne\-ma\-tik darum, herauszufinden,
welche Gelenkstellungen nötig sind, um eine bestimmte Zielposition zu erreichen.
Beide Verfahren sind essenziell für die Steuerung und Planung von Roboterbewegungen und stellen im
Vergleich zu den bisherigen Transformationen im Projekt eine sinnvolle und etwas komplexere Erweiterung dar.

Ein wichtiges Ziel der Bachelorarbeit ist es außerdem, das Framework auf dem JupyterLab-Server der Fachhochschule Dortmund bereitzustellen,
sodass es für Studierende leicht zugänglich wird. Dadurch soll sichergestellt werden, dass die neuen Funktionen nicht nur theoretisch existieren,
sondern auch direkt im Lehrbetrieb eingesetzt werden können. Die Idee ist,
dass Studierende ohne zusätzliche Installationen direkt im Browser mit dem Framework arbeiten und damit verschiedene Robotermodelle und kinematische Aufgaben ausprobieren können.

Auch die bestehenden interaktiven Elemente sollen weiterentwickelt werden.
Beispielsweise könnten bei der Rückwärtskinematik verschiedene Lösungsansätze visualisiert und direkt verglichen werden.
So bleibt das Framework weiterhin intuitiv und visuell verständlich.

Die Bachelorarbeit zu diesem Thema ist eine gute Gelegenheit,
das Framework inhaltlich und funktional weiterzuentwickeln und gleichzeitig einen praktischen Beitrag zur Lehre an der Fachhochschule Dortmund zu leisten.